\documentclass[10pt,a4paper]{amsart}
\usepackage[margin=19mm]{geometry}
\usepackage{amsmath,amssymb,mathtools}
\usepackage[scr=rsfs,cal=euler]{mathalfa}
\usepackage{xfrac}
\usepackage[unicode,hidelinks,bookmarks=false]{hyperref}
\newcommand{\ie}{i.e.\ }
\newcommand{\cf}{cf.\ }
\newcommand{\resp}{respectively\ }
\newcommand{\Subsection}{\subsection}
\providecommand{\nobreakdash}{\nobreak\hbox{-}}
\setlength{\emergencystretch}{2em}
\multlinegap0pt
\usepackage{amssymb}
\usepackage{xcolor}
\usepackage{tikz}
\newtheorem{lemma}{Lemma}
\title{Projective limits in Euclidean quantum field theory, II: Abelian gauge theory}
\author{Svetoslav Zahariev}
\date{Source: arXiv:2510.21399v2 (CC BY 4.0)}
\begin{document}
\maketitle

\noindent\emph{Source and licence.} Projective limits in Euclidean quantum field theory, II: Abelian gauge theory. Version arXiv:2510.21399v2 (CC BY 4.0), distributed by arXiv under the Creative Commons Attribution 4.0 International licence, http://creativecommons.org/licenses/by/4.0/, as recorded in the arXiv licence metadata for this exact version and shown on the abstract page https://arxiv.org/abs/2510.21399v2. Full licence notice: \texttt{licenses/arxiv-2510.21399-CC-BY-4.0.txt}.\par\smallskip

\noindent\textit{Editorial scope.} Counted unit: the article's lemma prpjistll (source0.tex lines 625--631). The article's complete original proof of it is delivered with it (source0.tex lines 632--652, one \texttt{proof} environment of 21 lines). No mathematics is written, repaired or re-derived: the statement and the proof are the article's own lines, in the article's own notation, and no retained line is substituted or re-flowed.  Retained uncounted as the source-local context the proof needs: lines 621--623.  Nothing was omitted from any retained span, so the package carries no omission record.  No retained line contains a citation, so the wrapper carries no bibliography and no bibliography entry.  No retained line contains a figure environment, an image-inclusion command or drawing code, so no image is delivered and the package declares no asset dependency.  Editorial formatting only, in addition: an AMS \texttt{amsart} preamble that restates the article's own declarations and macros used by the retained lines, namely the environments \texttt{lemma} on separate counters, together with the standard packages those lines need.  The retained environments print in this document as Lemma 1 in order of appearance, whereas the article prints the same items with the numbering of its own full document.  The complete native source of the article as distributed in the arXiv e-print for this exact version is bundled unchanged as \texttt{sources/187718/source0.tex}.
\begin{equation}\label{areanorminn}
\langle c_1, c_2 \rangle_T=\sum_{p \in T^{(2)}} \frac{1}{A(p)}c_1(p)c_2(p),\quad c_1,c_2 \in C^2(T;\mathfrak{g}),
\end{equation}

\begin{lemma}
The collection
\begin{equation}\label{prpjistll}
	\left\{(\mathtt{Im}\hspace{1pt}d_{T,1}^{\hspace{1pt}G},\mathtt{Im}\hspace{1pt}d_{T,1}^{\hspace{1pt}\mathfrak g},J_T),( P^2_{T,T'},(P^2_{T,T'})_{\bullet})\right\}_{(K,T) \in \mathcal{Q}^d},
\end{equation} 
where $\mathtt{Im}\hspace{1pt}d_{T,1}^{\hspace{1pt}\mathfrak g}$ is endowed with the inner product (\ref{areanorminn}),    is a projective system in $\mathbf{LG}$.
\end{lemma}

\begin{proof}
The identity $(P^2_{T,T'})^*_{\bullet}\circ J_{T'} =J_{T}\circ (P^2_{T,T'})^*$  follows from the naturality of the inclusion of the dual of a compact connected Abelian Lie group into the dual of its Lie algebra with respect to Lie group homomorphisms. It remains to show that $(P^2_{T,T'})_{\bullet}$ is  contraction. To this end, let $(K,T), (K',T') \in \mathcal{Q}^d$ with $(K,T)\preceq (K',T')$ and let $p$ be a 2-simplex in $T$. Then, choosing orientations of $p$ and the 2-simplices in $T'$ contained in $p$, we can write
$$
((P^2_{T,T'})_{\bullet}c)(p)=\sum_{q\subset p}\varepsilon(q,p)c(q), \quad c \in C^2(T';\mathfrak{g}),
$$
where the sum is over the 2-simplices in $T'$ contained in $p$ and the signs $\varepsilon(q,p)=\pm 1$ are determined by the choice of orientations. The Cauchy-Schwartz inequality implies
$$
\left|\sum_{q\subset p}\varepsilon(q,p)c(p)\right|^2
\leq
\left(\sum_{q\subset p}A(q)\right)
\left(\sum_{q\subset p}\frac{|c(q)|^2}{A(q)}\right),
$$
hence, noticing that $\sum_{q\subset p}A(q)=A(p)$, one has 
\begin{equation}\label{fracserin}
\frac{|(P^2_{T,T'})_{\bullet}c)(p)|^2}{A(p)}
\leq
\sum_{q\subset p}\frac{|c(q)|^2}{A(q)}.
\end{equation}
Finally, summing (\ref{fracserin}) over all 2-simplices in in $T$, we obtain  
$$\|((P^2_{T,T'})_{\bullet}c)\|^2_{T} \leq \|c\|^2_{T'}.$$
\end{proof}

\end{document}
